% =============================================================================
%  Section 3 — Three round regimes + Pearson correlations
% =============================================================================
\section{Three round regimes}
\label{sec:split_pearson}

CS2 rounds are not equivalent. \textbf{Pistol rounds} (R1, R13) start
from an economy reset to \$800 per player: money and weapon variables
are nearly constant. \textbf{Post-pistol rounds} (R2, R14) display a
structural economic asymmetry caused by the pistol outcome (winner on
a \emph{full-buy}, loser on \emph{eco/force-buy}). \textbf{Normal
rounds} cover the rest of the game, where each team freely chooses
its buy depending on accumulated cash, and where all buy types (eco,
half-buy, force-buy, full-buy) coexist. Training a single model on
these three regimes would mix three very different distributions; we
therefore separate the dataset into three splits at join time
(Table~\ref{tab:splits}). All subsequent analysis (correlations, PCA,
regression) is carried out separately on each split.

\begin{table}[H]
\centering
\small
\begin{tabular}{lrrr}
\toprule
\textbf{Split} & \textbf{Rounds} & \textbf{Columns} & \textbf{CT win rate} \\
\midrule
\texttt{df\_pistol}       & $2{,}945$  & 70  & 50.5~\% \\
\texttt{df\_post\_pistol} & $2{,}956$  & 105 & 52.3~\% \\
\texttt{df\_normal}       & $26{,}322$ & 106 & 51.1~\% \\
\bottomrule
\end{tabular}
\caption{Final partition of rounds.}
\label{tab:splits}
\end{table}

\subsection{Pearson correlations with \texttt{round\_winner}}
\label{subsec:pearson}

Our target is \texttt{round\_winner}, coded $Y = 1$ if CT wins the
round and $Y = 0$ otherwise. To quickly assess the strength of the
linear link between each feature $X$ and $Y$, we use the Pearson
correlation coefficient, bounded in $[-1, +1]$:
\[
r_{X,Y} = \frac{\operatorname{cov}(X,Y)}{\sigma_X\,\sigma_Y}
        = \frac{\sum_{i=1}^{n}(x_i-\bar x)(y_i-\bar y)}
               {\sqrt{\sum_{i=1}^{n}(x_i-\bar x)^2}\,\sqrt{\sum_{i=1}^{n}(y_i-\bar y)^2}}.
\]
For each round we form the \emph{deltas}
$\Delta X = X^{\text{CT}} - X^{\text{T}}$ of the economic and HLTV
features, and we compute $r$ between each $\Delta X$ and $Y$ on each
of the three splits (Figure~\ref{fig:correlations}).

Profiles are very contrasted. In \textbf{pistol} rounds all
correlations stay below $|r| < 0.06$: no individual feature
discriminates strongly. In \textbf{post-pistol} rounds the economic
deltas (\texttt{equipment\_value}, \texttt{money\_total},
\texttt{utility\_value}, weapon counts) reach $|r| \approx 0.6$, a
direct signature of the asymmetry created by the pistol outcome. In
\textbf{normal} rounds the same deltas remain at the top but with a
more moderate amplitude ($|r| \approx 0.3$). Across all three regimes
we observe a clear separation between the two feature families:
parser-derived economic variables carry most of the signal while
HLTV statistics (rating, ADR, round-win~\%, \ldots) remain
systematically at a lower correlation level.

\begin{figure}[!t]
\centering
\includegraphics[width=\linewidth]{report_pearson_combined.png}
\caption{Pearson correlations between each CT~$-$~T delta and
\texttt{round\_winner}, by regime. Blue: positive correlation
(feature pushes toward a CT win); orange: negative.}
\label{fig:correlations}
\end{figure}

% =============================================================================
